\section{Problem}
\label{sec:problem}

\begin{quote}\itshape
A web-based DT platform must render arbitrary visual substrates --- 3D assets,
photographs and live video --- loaded from arbitrary URLs, driven by heterogeneous live
data streams, across unrelated application domains, without rebuilding the stack each
time a new domain or medium appears. No open-source package solves this: every candidate
solves one domain deeply or all domains shallowly, and almost none treat an image or a
video feed as a legitimate place to put simulation output.
\end{quote}

\subsection{Six difficulties}

\textbf{P1 Domain rebuild.} 3D libraries are domain-shaped --- a BIM viewer models
storeys and IFC property sets, a geospatial engine models tiles and terrain, a robotics
viewer models joints. Serving four domains means either one poor fit in three of them or
four integrations to maintain. \citet{rundel2023leveraging} report the same pressure at
city scale: operational city DTs with simulation (Virtual Singapore, Herrenberg) were
\emph{impracticable} for small municipalities --- Virtual Singapore alone exceeded
USD~73~million over five years --- and their answer was procedural generation from
existing public geodata rather than bespoke authoring. The lesson generalises: derive
the scene from data that already exists.

\textbf{P2 Anchoring.} Something must connect a signal to a location. TwinMaker binds
entity/component/property onto a GLB submodel; iTwin.js keys on Element~ID; Tandem
classifies assets explicitly. All three \emph{author} the link and persist it; runtime
inference from coordinates or names breaks on re-export. DTaaS has an advantage here:
ThingsBoard already holds a device and asset registry, so the anchor namespace need not
be invented.

\textbf{P3 Sustained update.} Applying an encoding once is trivial; applying it to
$10^4$--$10^6$ elements at 10--60~Hz without rebuilding the scene graph is not. This
constrains technique choice far more than aesthetics.

\textbf{P4 Heterogeneity.} \citet{ali2024spatial} name this as open research: combining
3D geometry, time series and map layers legibly, and visualising correlations across
space and time, ``needs more research''. Part of the problem is genuinely unsolved, not
merely unbuilt.

\textbf{P5 Simulation coupling.} The problem commercial platforms serve worst. Tandem,
TwinMaker and Samsara visualise \emph{measurement} at high polish; simulation lives in
another tool with another viewer. Yet the residual between predicted and measured state
is the high-value signal, because it represents a physics violation rather than a
statistical oddity. \citet{liao2025digital} co-visualise sensed and simulated
deformation of a soft robot in one AR scene, reporting errors of 0--7.14\%.

\textbf{P6 Substrate plurality.} Most libraries assume the answer is a 3D scene. In
practice a large share of requests are a \emph{photograph} (traffic still, room, plant
floor, scanned plan) on which simulation should be drawn; a \emph{live video feed} to be
shown beside the simulation; or a \emph{rendered output file} from a workspace execution
that needs only placing and time-aligning. Treating these as second-class leaves the
visualisation layer unusable until someone commissions a 3D model.

\begin{keypoint}
DTaaS is unusually placed on P5: its INTO-CPS lineage means FMI co-simulation already
produces the predicted channel every surveyed commercial platform lacks, and InfluxDB
already provides somewhere to put both channels. The cost of exploiting this is small
\emph{provided} measured and simulated state are two channels of one signal from the
start.
\end{keypoint}

\subsection{Why the obvious approaches fail}

\textbf{One 3D library} fails on P1 and P6. Pascal (\pkg{@pascal-app/*}, MIT) gives a
full architectural editor but its document model is sites\,$\rightarrow$\,buildings\,%
$\rightarrow$\,levels\,$\rightarrow$\,walls, in which a pump, a fjord or a photograph is
not a citizen, and its viewer exposes no live-data binding. \pkg{@thatopen/components}
has the best data semantics of any BIM option and is likewise BIM-shaped. CesiumJS is
decisive at city scale and clumsy for one machine. \textbf{Everything on three.js} fails
on effort, and is the wrong tool for annotating a JPEG. \textbf{A game engine} works ---
\citet{rundel2023leveraging} use Unity+SUMO via TraCI, \citet{vasilijevic2024trondheim}
Unity+Cesium --- but produces an artefact hosted and versioned outside the platform; its
proper home in DTaaS is the workspace desktop (\S\ref{sec:workspace}).
\textbf{Grafana and ThingsBoard alone} solve charting completely and the spatial problem
not at all. \textbf{A commercial platform} does not solve P5, and only iTwin.js (MIT) is
installable.

\subsection{Requirements}

\begin{longtable}{@{}L{1.1cm}L{13.4cm}@{}}
\toprule
\textbf{ID} & \textbf{Requirement} \\
\midrule
\endhead
\req{1} & \textbf{Substrate independence.} Adding a domain or a medium means adding an
adapter behind a stable interface, not forking the package. \\
\req{2} & \textbf{Authored anchoring.} Signal-to-location bindings are explicit,
persisted, versionable and human-editable. \\
\req{3} & \textbf{Transport agnosticism.} MQTT and AMQP (both via the installed
RabbitMQ), ThingsBoard telemetry, InfluxDB, SQL, HTTP and files normalise to one stream
shape. \\
\req{4} & \textbf{Temporal addressing.} Every encoding reads state at a playhead, never
``latest value''. \\
\req{5} & \textbf{Dual-channel state.} Measured and simulated values of one signal are
first-class and separately addressable. \\
\req{6} & \textbf{Sustained update.} $O(\text{changed elements})$ per frame, no
scene-graph rebuild. \\
\req{7} & \textbf{Composability as a DTaaS asset.} A visualisation is storable in the
library, selectable in composition, shareable between users. \\
\req{8} & \textbf{Licence cleanliness.} MIT, Apache-2.0 or BSD only; network copyleft is
disqualifying for a platform client. \\
\req{9} & \textbf{Graceful degradation.} A twin with no 3D asset still gets a usable
view; a browser without WebGPU still renders. \\
\req{10} & \textbf{Reuse of the installed stack.} Consume InfluxDB, RabbitMQ,
ThingsBoard, Grafana and GitLab as deployed; no parallel broker, historian, dashboard or
file store. \\
\req{11} & \textbf{Media synchronisation.} Where video and simulation are shown
together, their time alignment is explicit, measurable and correctable. \\
\bottomrule
\end{longtable}

\noindent
\req{8} eliminates \pkg{xeokit-sdk} and \pkg{realvirtual WEB} (both AGPL-3.0) and flags
\pkg{Potree} and \pkg{loaders.gl} (\texttt{NOASSERTION}) for checking. \req{10} is what
most changes the plan relative to a greenfield design.

